%%======%%======%0%======%%=======%%
\vspace{5mm}
\section{Limitação}
%%======%%======%0%======%%=======%%

Apesar de o sistema desenvolvido atender à proposta inicial de apresentar o Laboratório de Investigação de Filmes Semicondutores e organizar suas principais informações, algumas limitações ainda permanecem na versão atual. Isso ocorre principalmente porque o projeto foi desenvolvido dentro do escopo da primeira etapa da disciplina SCOM, que possui como foco a construção e avaliação da interface web. Dessa forma, algumas funcionalidades foram estruturadas apenas até o nível necessário para demonstrar seu funcionamento no protótipo.

Uma das principais limitações está relacionada à persistência dos dados. Parte das informações utilizadas pelo sistema é armazenada por meio do \textit{LocalStorage} do navegador. Essa abordagem é suficiente para o funcionamento de uma aplicação local e para demonstrar o fluxo de cadastro e consulta, porém não representa uma solução adequada para uma utilização simultânea por vários usuários. Os dados permanecem associados ao navegador utilizado e não existe, na implementação atual, uma base de dados central responsável por armazenar e sincronizar as informações.

Essa limitação aparece principalmente no repositório de trabalhos e no gerenciamento das informações apresentadas nas páginas. O sistema já possui uma estrutura para cadastro de trabalhos e organização dos dados, mas ainda não existe uma camada de servidor responsável por validar, armazenar e disponibilizar essas informações de maneira centralizada. Portanto, o cadastro realizado no protótipo não deve ser entendido como um sistema definitivo de gerenciamento da produção científica do laboratório.

A autenticação também possui uma limitação semelhante. Embora exista uma página de login e uma lógica de controle de sessão no lado do cliente, ela ainda não corresponde a um mecanismo de autenticação de usuários adequado para uma aplicação real. Como a aplicação é predominantemente estática e executada no navegador, não existe atualmente um servidor responsável por verificar credenciais e controlar permissões de maneira segura. Assim, essa funcionalidade deve ser considerada parte da estrutura inicial do sistema, e não um controle de acesso definitivo.

Outra limitação está no agendamento dos equipamentos. A página de máquinas já apresenta os equipamentos disponíveis e a interface foi preparada considerando a possibilidade de solicitação de uso, porém o controle efetivo de disponibilidade ainda não está implementado de forma centralizada. Não há, por exemplo, um calendário compartilhado que impeça dois usuários de solicitarem o mesmo equipamento para o mesmo período. Essa parte foi deixada como uma evolução natural do projeto, principalmente porque exigiria uma estrutura de dados persistente e regras para tratamento dos horários.

Também existe uma limitação relacionada aos próprios documentos disponibilizados no repositório. Nem todos os trabalhos científicos podem simplesmente ter seus arquivos PDF disponibilizados publicamente, devido às condições de publicação e aos direitos associados a determinados periódicos. Por esse motivo, uma versão futura do sistema deverá diferenciar a existência do registro bibliográfico da disponibilização do arquivo completo. Em alguns casos, pode ser mais apropriado apresentar os dados da publicação e direcionar o usuário para a fonte oficial.
